% M11, 11.09.2026: Gestaltungsfolgen bündeln; Mechanik bleibt in 5.2--5.9.
% Governance-Reichweite nach 5.6; Initialisierung und F1 nicht verschmolzen.

\section{Architektursynthese und Gestaltungsfolgen}
\label{sec:architektursynthese-gestaltungsfolgen}

Die Architektur verbindet quellengebundene Verarbeitung mit einem
fortschreibbaren Projektbestand. Dabei werden fachliche Vorschläge,
übernommene Inhalte und operative Sichten unterscheidbar gehalten.
Der Durchgang in Abschnitt~\ref{sec:endtoend-pfade} zeigt am
Besuchsbeispiel, wie diese Trennung die einzelnen Übergänge prüfbar
macht; die dort ebenfalls erläuterte Initialisierung bleibt ein
eigener Vorgang.

Aus der Aufgabenverteilung folgt ein begrenzter, aber bewusst
erhaltener Handlungsspielraum für Agenten. Der Steward wählt
Fähigkeiten für die Interaktion aus; der Core-Analyst untersucht den
Bestand auf mögliche Lücken und Widersprüche. Fachliche Änderungen
führen beide über die vorgesehenen Verarbeitungs- und Governancepfade
herbei. Typisierte Übergaben, feste Ablaufregeln und Prüfungen
begrenzen dabei die möglichen Zustandswirkungen. Sie garantieren
keine semantisch richtigen oder bei Wiederholung identischen
Modellergebnisse.

Die kontrollierte Fortschreibung bindet zugleich menschliche
Entscheidungsarbeit ein. Die Autorisierung erfolgt nach den Regeln
des jeweiligen fachlichen Pfads; die gemeinsame Speicheroperation
prüft anschließend die strukturelle Zulässigkeit. Beide Kontrollen
ergänzen sich, haben aber unterschiedliche Aufgaben. Der Erst-Seed
ohne Einzelautorisierung bleibt die in
Abschnitt~\ref{sec:governance-mutation} ausgewiesene Einschränkung.
Offene Fragen und Ablehnungen bleiben im Projektbestand erhalten,
können jedoch weiterhin menschliche Klärung erfordern. Ihre
Speicherung löst nicht automatisch jede weitere Verarbeitung aus.

Auch die Außenkopplung besitzt eine klare Richtung: GitHub und
Dokumente stellen abgeleitete Sichten dar. Externe Rückmeldungen
gelangen über den vorgesehenen GitHub-Rückweg erneut in die
kontrollierte Verarbeitung; für Dokumentsichten besteht kein
entsprechender Rückweg. Ob die Verarbeitung Inhalte ausreichend
erhält und die Beziehungen tatsächlich nachvollziehbar bleiben,
ist damit noch nicht beantwortet. Ausgewählte Eigenschaften werden
in Kapitel~\ref{chap:evaluation} untersucht und in
Kapitel~\ref{chap:diskussion-fazit} eingeordnet.

Kapitel~\ref{chap:maf-realisierung} zeigt zunächst, welche Teile
dieser Architektur das Microsoft Agent Framework trägt und welche
Regeln sowie Sicherungen das Projekt ergänzt.
